\section{Question-by-question coverage}\label{sec:coverage}

Question identifiers preserve the proposal's research-area numbering and the
order of questions within each area. R8 and R9 are both in its
\code{sections/12-ranking.tex}. The consolidated expert-contact section
reorders these themes; it does not add a separate set of research questions.
The summaries below paraphrase the questions rather than reproduce them.

\begin{longtable}{@{}L{17mm}L{43mm}L{82mm}r@{}}
\caption{All 31 expert questions and the answer supplied.}\\
\toprule
ID & Question theme & Main answer & Page\\
\midrule
\endfirsthead
\toprule ID & Question theme & Main answer & Page\\\midrule
\endhead
R1.Q1 & Shared metavariables & Conditional factorization, persistent branches, interface-aware tabling. & \pageref{q:R1.Q1}\\[4pt]
R1.Q2 & Admissible remaining cost & Maximum of valid bounds; sums only with cost separability. & \pageref{q:R1.Q2}\\[4pt]
R1.Q3 & Online learning and ranking & Freeze the objective; adapt priorities with recorded logical epochs. & \pageref{q:R1.Q3}\\[4pt]
R1.Q4 & Deadline determinism & Identical timed answer sets are generally impossible; replay prefixes. & \pageref{q:R1.Q4}\\[4pt]
R2.Q1 & Minimal carrier type & Semantic quotient and finite-state lower bounds; bounded syntax search. & \pageref{q:R2.Q1}\\[4pt]
R2.Q2 & Polymorphic alphabet & Naturality yields position tests under explicit assumptions. & \pageref{q:R2.Q2}\\[4pt]
R2.Q3 & Backward fusion & Use lifting and auxiliary-function precedents with invariant obligations. & \pageref{q:R2.Q3}\\[4pt]
R2.Q4 & Recursor motive basis & Bounded typed carrier/motive grammar, including parameter binders. & \pageref{q:R2.Q4}\\[4pt]
R2.Q5 & Generalization heuristics & Generalize changing state and synthesize strengthening candidates. & \pageref{q:R2.Q5}\\[4pt]
R3.Q1 & Higher-order/live evaluation & Smyth already supports higher-order examples; dependent fibers add work. & \pageref{q:R3.Q1}\\[4pt]
R3.Q2 & Incremental normalization & Dependency read sets with rollback-safe assignment identities. & \pageref{q:R3.Q2}\\[4pt]
R3.Q3 & Conservative fast evaluator & A typed fragment with checked evaluation/refutation evidence. & \pageref{q:R3.Q3}\\[4pt]
R3.Q4 & Top-down angelic guarantees & Finite bounded completeness under decisive constraints and fairness. & \pageref{q:R3.Q4}\\[4pt]
R4.Q1 & Motive enumeration & Finite abstraction templates plus bounded accumulator closure. & \pageref{q:R4.Q1}\\[4pt]
R4.Q2 & Canonical transports & Normalize indices with evidence; do not erase required casts. & \pageref{q:R4.Q2}\\[4pt]
R4.Q3 & Unfolding policy & Staged head exposure informed by Canonical and documented sauto options. & \pageref{q:R4.Q3}\\[4pt]
R4.Q4 & Lean motive API & Reuse prefix construction and induction with checked preconditions. & \pageref{q:R4.Q4}\\[4pt]
R5.Q1 & Interpreter versus realization & Propositional simulation and equation proofs; fuel exhaustion is unknown. & \pageref{q:R5.Q1}\\[4pt]
R5.Q2 & Top-down angelic composition & Shared typed call constraints and explicit assumption discharge. & \pageref{q:R5.Q2}\\[4pt]
R5.Q3 & Small recursion basis & Paramorphism-oriented templates; measure coverage by capability. & \pageref{q:R5.Q3}\\[4pt]
R6.Q1 & Value/proof interleaving & Dependency-aware readiness and local invariant verification conditions. & \pageref{q:R6.Q1}\\[4pt]
R6.Q2 & Recognizable induction class & Sufficient syntactic residual-theory tests, not universal recognition. & \pageref{q:R6.Q2}\\[4pt]
R6.Q3 & Usable induction automation & Core meta APIs and functional induction, with bounded branch solvers. & \pageref{q:R6.Q3}\\[4pt]
R6.Q4 & Proof-debt scheduling & Deterministic fair quanta plus cost/benefit estimates. & \pageref{q:R6.Q4}\\[4pt]
R7.Q1 & Dependent type abstraction & Constraint-carrying hyperedges with a proved over-approximation. & \pageref{q:R7.Q1}\\[4pt]
R7.Q2 & Classes in refinement nets & Hoogle+ dictionary passing; preserve Lean's contextual dictionaries. & \pageref{q:R7.Q2}\\[4pt]
R7.Q3 & Program relevance learning & Data-goal lookup exists; build example-aware retrieval with clean splits. & \pageref{q:R7.Q3}\\[4pt]
R8.Q1 & Ranking objective & Explicit description-length prior and evaluated user-facing preferences. & \pageref{q:R8.Q1}\\[4pt]
R8.Q2 & Candidate equivalence & Separate definitional/proved equality from provisional probe buckets. & \pageref{q:R8.Q2}\\[4pt]
R9.Q1 & Model carrier accuracy & No directly transferable number; measure four separate success rates. & \pageref{q:R9.Q1}\\[4pt]
R9.Q2 & Deterministic auditability & Separate proof, search, and model replay; record restricted proposals. & \pageref{q:R9.Q2}\\
\bottomrule
\end{longtable}

\section{Source audit and scope of verification}\label{sec:audit}

The repository reading was pinned to commit
\code{784664ff34e819422510a4d470ab07fe48bb736b}, dated September 18, 2026.
All nine research areas and their questions were read in their source TeX,
together with the evidence, engineering, adaptation, and consolidated expert
sections, the bibliography, the root README, and the implementation notes.
The analysis of current code uses selected modules rather than a claim to
have audited the entire repository. Stable source URLs are in the bibliography.

\begin{longtable}{@{}L{65mm}L{88mm}@{}}
\toprule
Source & What was checked\\
\midrule
\endhead
\code{Leant2/Search/Core.lean}, lines 1--440 & Provider indexing, local frontier, snapshots, delayed substitution, residual decisions and blockers, and the closed-goal proof portfolio.\\[4pt]
\code{Leant2/Native/Transaction.lean} & Search-context references, work accounting, deadlines, and rollback behavior.\\[4pt]
\code{Leant2/Accept/Gate.lean} & Joint universe generalization, synchronous checking, supplied proof type, and axiom audit.\\[4pt]
Lean 4.34.0: \code{src/Lean/Meta/Tactic/Induction.lean} & Index preconditions, motive prefix, dependency reversion, and returned branch substitutions.\\[4pt]
Lean 4.34.0: \path|src/Lean/Meta/Tactic/LibrarySearch.lean|, lines 1--80 & Data-goal example and direct \code{Grind.main} adapter.\\[4pt]
Official Lean manual, current version shown as 4.35.0-rc2 & Functional induction, recursive definitions, and proof-validation guidance; not treated as pinned 4.34.0 source.\\
\bottomrule
\end{longtable}

Primary papers were used for the claims about existing synthesis methods.
Where a current work is a preprint, it is identified as such; its results are
not promoted to a theorem about Leant2. Bibliographic details in the original
proposal were not simply copied: the investigation specifically checked the
features relevant to each question. Absence of a directly matching benchmark
in the sources examined does not prove that no such experiment exists.

No Lean executable was available in the working environment, and the Leant2
baseline and proposed adapters were not run. The only executable evidence
newly produced here is the finite Python suite described in
Section~\ref{sec:evaluation}. Ordinary mathematical proofs and implementation
recommendations must not be represented as a completed Lean formalization.

\section{Using the accompanying package}

Build the article from the package root with
\begin{lstlisting}
latexmk -pdf -interaction=nonstopmode -halt-on-error article.tex
\end{lstlisting}
A standard TeX distribution with the listed packages is sufficient; no
bibliography service, external figure download, or font file is required.
The bibliography is included as \code{references.tex}. Without \code{latexmk},
run \code{pdflatex} three times so the contents and references settle.

Run the finite-model checks with Python 3.9 or newer:
\begin{lstlisting}
python checks/check_models.py --output checks/results.json
\end{lstlisting}
Use a normal Python invocation without \code{-O}, because the checks use
assertions. The script needs no third-party package. Its JSON output records
which cases were examined and the discovered finite-state witness.

The package also includes a README and a provenance manifest. The manifest
records the repository commit and the distinction between inspected code,
published results, written mathematical arguments, and executed finite checks.
It is a provenance aid, not a cryptographic certificate of the article's
mathematical conclusions.
